\documentclass[11pt,a4paper]{article}
\pdfoutput=1

\usepackage{jheppub}
\usepackage{amsmath}
\usepackage{amssymb}
\usepackage{amsthm}
\newtheorem{proposition}{Proposition}
\usepackage{booktabs}
\usepackage{graphicx}
\usepackage{hyperref}
\usepackage{microtype}
\sloppy

% Loud marker for any number whose artifact was missing at build time.  A
% missing input must be visible in the PDF, never a silently stale value.
\newcommand{\ARTIFACTMISSING}[1]{\textbf{\color{red}[ARTIFACT MISSING: #1]}}
\usepackage{xcolor}

% GENERATED FILE -- do not edit.
% Produced by manuscript/scripts/make_numbers.py from result artifacts.

\newcommand{\dimAten}{14}
\newcommand{\dimBten}{12}
\newcommand{\dimGten}{12}
\newcommand{\dimPten}{2}
\newcommand{\dimDten}{11}
\newcommand{\dimQten}{3}
\newcommand{\capBG}{10}
\newcommand{\capPB}{1}
\newcommand{\capPG}{0}
\newcommand{\capDB}{9}
\newcommand{\sumPB}{13}
\newcommand{\nIncidencePrimes}{2}
\newcommand{\dimBcapP}{1}
\newcommand{\dimBplusP}{13}
\newcommand{\publishedPrimitiveContent}{11}
\newcommand{\bridgePrimes}{32719, 32749}
\newcommand{\nBridgePrimes}{2}
\newcommand{\starSquared}{1}
\newcommand{\dimSelfdual}{126}
\newcommand{\dimAntiSelfdual}{126}
\newcommand{\bridgeForwardRank}{126}
\newcommand{\dimGammaTraceless}{126}
\newcommand{\gammaTraceRank}{10}
\newcommand{\comparisonPrime}{32749}
\newcommand{\nCommonSamples}{78}
\newcommand{\nSampleSparse}{10}
\newcommand{\nSampleStructured}{8}
\newcommand{\nSampleGeneric}{40}
\newcommand{\nSampleHoldout}{20}
\newcommand{\traceRankDegFour}{1}
\newcommand{\spinorRankDegFour}{1}
\newcommand{\spansEqualDegFour}{yes}
\newcommand{\holdoutDegFour}{yes}
\newcommand{\traceRankDegSix}{2}
\newcommand{\spinorRankDegSix}{2}
\newcommand{\spansEqualDegSix}{yes}
\newcommand{\holdoutDegSix}{yes}
\newcommand{\traceRankDegEight}{7}
\newcommand{\spinorRankDegEight}{6}
\newcommand{\spansEqualDegEight}{no}
\newcommand{\holdoutDegEight}{no}
\newcommand{\traceRankDegTen}{14}
\newcommand{\spinorRankDegTen}{14}
\newcommand{\spansEqualDegTen}{yes}
\newcommand{\holdoutDegTen}{yes}
\newcommand{\archiveJacobianRank}{59}
\newcommand{\archiveJacobianRunsAgree}{yes}
\newcommand{\nArchiveCandidates}{83}
\newcommand{\nArchivePortGraphs}{70}
\newcommand{\nArchiveStructured}{13}
\newcommand{\nArchiveJacobianRuns}{3}
\newcommand{\nFloatRuns}{2}
\newcommand{\nFloatNondegenerate}{1}
\newcommand{\nFloatDegenerate}{1}
\newcommand{\nFloatInconclusive}{0}
\newcommand{\floatNondegenerateRanks}{81}
\newcommand{\exactJacRank}{81}
\newcommand{\exactJacRunsAgree}{yes}
\newcommand{\exactJacPoints}{2}
\newcommand{\exactJacPointWord}{points}
\newcommand{\exactJacPrimeWord}{prime}
\newcommand{\exactJacPrimes}{32749}
\newcommand{\nExactJacPrimes}{1}
\newcommand{\exactJacRows}{83}
\newcommand{\exactJacCols}{126}
\newcommand{\exactJacScheduled}{83}
\newcommand{\exactJacEvaluated}{83}
\newcommand{\exactJacErrors}{0}
\newcommand{\exactJacZeroRows}{0}
\newcommand{\exactJacScheduleComplete}{yes}
\newcommand{\exactJacEulerPass}{yes}
\newcommand{\charZeroLowerBound}{81}
\newcommand{\cumRankDegFour}{1}
\newcommand{\cumRankDegSix}{3}
\newcommand{\cumRankDegEight}{9}
\newcommand{\cumRankDegTen}{21}
\newcommand{\cumRankDegTwelve}{81}
\newcommand{\minorSize}{81}
\newcommand{\minorPrimes}{32749}
\newcommand{\nMinorPrimes}{1}
\newcommand{\minorDetNonzero}{yes}
\newcommand{\dEightTraceRank}{7}
\newcommand{\dEightSpinorRank}{7}
\newcommand{\dEightUnionRank}{7}
\newcommand{\dEightSpinorRowsMin}{21}
\newcommand{\dEightSpinorRowsMax}{23}
\newcommand{\dEightPrimes}{4}
\newcommand{\dEightPrimesEqual}{4}
\newcommand{\dEightIndispensable}{tensor word}
\newcommand{\dEightRankWithoutWords}{6}
\newcommand{\nTraceTests}{207}
\newcommand{\nBridgeTests}{72}
% analytic constants -- literature or elementary, not computed here
\newcommand{\dimFiveFormComponents}{252}
\newcommand{\dimSelfDualModule}{126}
\newcommand{\dimLorentzAlgebra}{45}
\newcommand{\genericFunctionalDim}{81}
\newcommand{\dimSymSpinor}{136}
\newcommand{\spinorDim}{16}
\newcommand{\dimQtenMinusOne}{2}


\newcommand{\HUMAN}[1]{\textbf{[#1 --- HUMAN ACTION REQUIRED]}}

\title{An exact degree-ten classification of local invariants of the
       ten-dimensional self-dual five-form}

% Placeholders below are deliberately plain text: they must compile, and they
% must be impossible to miss in the PDF.  See
% submission_candidate/AUTHORSHIP_DECISION_REQUIRED.md.
\author[a]{AUTHOR LIST -- HUMAN ACTION REQUIRED}
\affiliation[a]{AFFILIATIONS -- HUMAN ACTION REQUIRED}
\emailAdd{corresponding.author@placeholder.invalid}

\abstract{%
The self-dual five-form field strength in ten dimensions carries a large ring of
Lorentz invariants, of which only a handful have been written down explicitly.
We determine the degree-ten graded piece of that ring exactly. Working over
finite fields with certified linear algebra, we construct a fourteen-dimensional
atlas of degree-ten invariants and resolve its internal structure completely:
the two-dimensional subspace of products of lower-degree invariants, the
twelve-dimensional span of the published candidate structures, the
twelve-dimensional span of the graph generators, and the eleven-dimensional
subspace reachable by the stress-tensor flow that motivates these theories. The
resulting quotient is three-dimensional --- exactly so over the rationals, not
only modulo a prime --- and we give explicit compact representatives for it. We show that two decompositions which both read twelve
plus two are not the same decomposition: the published span contains one product
direction, so its non-product content is eleven. Independently, we implement a
bridge between the tensor description and a spinor-variable description, fixing
every convention explicitly and verifying the map is an isomorphism onto the
gamma-traceless symmetric chiral square. The two descriptions are then evaluated
at a common registry of five-forms and their invariant spans compared directly.
We also replace the floating-point Jacobian rank computation used previously by
an exact modular one over the complete candidate selection. Because the
coordinate basis is integral, the resulting rank is an unconditional lower bound
in characteristic zero rather than a tolerance-dependent estimate, and it is
witnessed by an explicit non-vanishing minor; it meets the analytic count of
eighty-one from below.%
}

\keywords{Extended Supersymmetry, p-branes, Differential and Algebraic Geometry}

\begin{document}
\maketitle

\begin{center}
\fbox{\parbox{0.9\textwidth}{\small\textbf{Draft status.}
This is a private submission candidate. The author list, affiliations,
acknowledgments, funding, licence and data DOI are unresolved placeholders and
are listed in \texttt{submission\_candidate/AUTHORSHIP\_DECISION\_REQUIRED.md}.
The physics interpretation of section~\ref{sec:physics} and the novelty
statements require coauthor confirmation. Not for distribution.}}
\end{center}

%=============================================================================
\section{Introduction}
\label{sec:intro}

A self-dual five-form field strength $F_{\mu_1\ldots\mu_5}$ in ten Lorentzian
dimensions has $\dimSelfDualModule$ independent components. The Lorentz group
$SO(1,9)$ has dimension $\dimLorentzAlgebra$ and acts with trivial generic
stabiliser, so a generic orbit is $\dimLorentzAlgebra$-dimensional and the number
of functionally independent invariants is~\cite{Hutomo:2025chiral}
\begin{equation}
  \dimSelfDualModule - \dimLorentzAlgebra = \genericFunctionalDim .
  \label{eq:eightyone}
\end{equation}
This count is analytic and is not in question~\cite{Hutomo:2025chiral}. What is much less accessible is
the \emph{ring} of invariants: which polynomial invariants exist at each degree,
how they are related, and which of them can be written down.

The contrast with lower dimensions is stark. In four and six dimensions the
analogous problem has a handful of generators and the invariants relevant to
non-linear duality-invariant theories can be listed on one line. In ten
dimensions the number in~\eqref{eq:eightyone} means that no such listing is
possible, and the explicit structures available in the literature are a small
and unsystematic sample of a much larger space.

\paragraph{What was previously available.} Twelve explicit degree-ten
candidate structures and three degree-twelve structures appear in the source
literature~\cite{Cederwall:2025invariants}. They are given as index expressions, they are not claimed to be a
basis, and their linear relations were not determined.

\paragraph{The gap addressed here.} Three questions were open at degree ten.
First, what is the actual dimension of the degree-ten graded piece, and do the
published structures span it? Second, the theories that motivate this problem
are generated by a stress-tensor flow; which degree-ten invariants does that flow
actually reach, and what is the quotient? Third, the same invariant ring can be
presented in spinor variables, and a separate implementation exists that does so
--- but the two implementations use different frames, and no map between them had
been constructed, so their agreement had never been tested at the level of
values.

\paragraph{Results.} We answer all three, exactly.

\begin{enumerate}
\item The degree-ten atlas $A_{10}$ has dimension $\dimAten$
      (section~\ref{sec:atlas}).
\item Its internal structure is determined completely
      (section~\ref{sec:decomp}): the product subspace $P_{10}$ has dimension
      $\dimPten$, the published span $B_{10}$ has dimension $\dimBten$, the graph
      generator span $G_{10}$ has dimension $\dimGten$, the stress-flow reachable
      subspace $D_{10}$ has dimension $\dimDten$, and the quotient
      $Q_{10} = A_{10}/D_{10}$ has dimension $\dimQten$. All pairwise
      intersections and sums are certified.
\item The decomposition $A_{10} = G_{10}\oplus P_{10}$ holds, but the published
      span is \emph{not} a complement of the products: $\dim(B_{10}\cap P_{10})
      = \dimBcapP$, so the non-product content of the published span is
      $\publishedPrimitiveContent$, not $\dimBten$.
\item A convention-controlled bridge to the spinor description is implemented and
      verified to be an isomorphism onto the gamma-traceless
      $\dimGammaTraceless$ (section~\ref{sec:bridge}).
\item The floating-point Jacobian used previously is replaced by an exact modular
      one (section~\ref{sec:jacobian}), which converts numerical evidence into a
      rigorous bound.
\end{enumerate}

\paragraph{What remains unclassified.} Degree twelve is background here, not a
result. We do not present a generating set for the full invariant ring, we do not
claim an all-orders statement, and the reverse search of
section~\ref{sec:reverse} is bounded and is not an exhaustive enumeration. These
limits are stated precisely in section~\ref{sec:limitations}.

%=============================================================================
\section{The self-dual five-form and its invariants}
\label{sec:setup}

\subsection{Conventions}

We work with metric signature $(-,+,\ldots,+)$ and
$\epsilon_{0 1 \ldots 9} = +1$. The Hodge dual of a $p$-form is
\begin{equation}
  (\star F)_{I} = \tfrac{1}{(d-p)!}\,\epsilon_{I}{}^{J} F_{J},
\end{equation}
and for $d=10$, $p=5$,
\begin{equation}
  \star^2 = (-1)^{p(d-p)}\,\mathrm{sign}(\det\eta)
          = (-1)^{25}\cdot(-1) = +\starSquared ,
\end{equation}
so real self-dual five-forms exist. Self-duality means $\star F = +F$ and the
projector is $(1+\star)/2$. The $\dimFiveFormComponents$ components of a general
five-form split as $\dimSelfdual + \dimAntiSelfdual$. Every one of these
statements is verified at runtime rather than assumed; see
appendix~\ref{app:conventions}.

\subsection{Invariants, and five things called ``rank''}
\label{sec:dictionary}

A recurring source of confusion is that several distinct numbers are all
naturally called a rank or a dimension. We fix the vocabulary once, and use it
without exception.

\begin{description}
\item[Atlas dimension.] The dimension of the space of polynomial invariants of a
  given $F$-degree. Graded, not cumulative. Exact.
\item[Evaluation rank.] The rank of the matrix of invariant values over a finite
  set of sample five-forms. Equals the atlas dimension once the sample set is
  rich enough, and is what a computation actually measures.
\item[Jacobian rank.] The rank of $\partial I_i/\partial c_r$ at a point. This
  counts \emph{functionally} independent invariants and is cumulative over all
  degrees. It is bounded by~\eqref{eq:eightyone}.
\item[Hilbert-series dimension.] The graded dimension predicted by a generating
  function. A count of generators, not of functions.
\item[Quotient dimension.] The dimension of $A_{10}/D_{10}$, a quotient by the
  stress-flow reachable subspace --- unrelated to any product decomposition.
\end{description}

These genuinely differ, and the cleanest demonstration is at the top of the
range rather than in the middle. The candidate family whose Jacobian we compute
in section~\ref{sec:cert81} has $\exactJacRows$ members and Jacobian rank
$\exactJacRank$, so exactly two functional relations hold among functions that
are pairwise distinct as polynomials. Both are degree-twelve phenomena: the
cumulative Jacobian ranks are $\cumRankDegFour, \cumRankDegSix,
\cumRankDegEight, \cumRankDegTen, \cumRankDegTwelve$ against per-degree candidate
counts that each contribute in full up to degree ten.

A caution that cost us a wrong sentence in an earlier draft: one may not subtract
a Jacobian rank taken over the spinor candidate family from a cumulative atlas
dimension taken on the tensor side. The two are indexed by different sets --- at
degree eight the tensor atlas has $\traceRankDegEight$ dimensions while the
spinor family contributes $\spinorRankDegEight$ candidates --- and their
difference is not a count of anything. Table~\ref{tab:dims} collects the values
with each one's space named.

\begin{table}[t]
\centering
\begin{tabular}{lrll}
\toprule
space & dimension & arithmetic & grading \\
\midrule
$\dim$ self-dual module & 126 & analytic & -- \\
$\dim \mathfrak{so}(1,9)$ & 45 & analytic & -- \\
generic functional dimension & 81 & analytic & cumulative \\
atlas $A_{10}$ & 14 & exact modular & graded, degree 10 \\
published span $B_{10}$ & 12 & exact modular & graded, degree 10 \\
graph generators $G_{10}$ & 12 & exact modular & graded, degree 10 \\
products $P_{10}$ & 2 & exact modular & graded, degree 10 \\
reachable $D_{10}$ & 11 & exact modular & graded, degree 10 \\
quotient $Q_{10}$ & 3 & exact modular & graded, degree 10 \\
\bottomrule
\end{tabular}
\caption{Degree-ten dimensions. Each row dimensions a different space; several coincide numerically and must not be conflated. Modular values are identical at all 2 primes tested.}
\label{tab:dims}
\end{table}


%=============================================================================
\section{Exact computational framework}
\label{sec:framework}

\subsection{Graphs, canonicalisation and modular evaluation}

An invariant is represented as a contraction graph: nodes are copies of $F$,
edges are index contractions carrying exactly one factor of the metric. Graphs
are canonicalised so that two labellings of the same contraction are identified.
Evaluation is carried out over $\mathbb{F}_p$ at
$p \in \{\bridgePrimes\}$ rather than in floating point. The reason is
structural rather than one of convenience: invariant values span an enormous
dynamic range at high degree, so a floating-point rank requires choosing a
tolerance, whereas over $\mathbb{F}_p$ vanishing is unambiguous.

\begin{figure}[t]
\centering
\includegraphics[width=0.95\textwidth]{figures/trace_pipeline.pdf}
\caption{The tensor-side pipeline. Nothing downstream of the rank sieve uses a
tolerance, because nothing upstream of it uses floating point.}
\label{fig:pipeline}
\end{figure}

\subsection{What modular arithmetic does and does not establish}
\label{sec:modular-caveat}

Rank over $\mathbb{F}_p$ at a random specialisation is a \emph{lower} bound on
the rank in characteristic zero, unconditionally, provided the matrix is the
reduction of an integer matrix --- reduction can only drop rank. It can fail to
reach the characteristic-zero rank only on a proper subvariety, whose measure is
$O(\deg/p)$ by the Schwartz--Zippel bound~\cite{Schwartz:1980,Zippel:1979}. We therefore use agreement at two
primes as strong evidence, and never as a proof, for any statement that requires
the \emph{upper} direction. Where an upper bound is needed we take it from an
analytic argument, as in~\eqref{eq:eightyone}.

Sampling and validation are separated: bases are fitted on one set of five-forms
and every claimed relation is then checked on a disjoint holdout set that took no
part in the fit.

%=============================================================================
\section{The degree-ten atlas}
\label{sec:atlas}

Enumerating degree-ten contraction graphs, canonicalising, and computing the
evaluation rank over $\mathbb{F}_p$ gives
\begin{equation}
  \dim A_{10} = \dimAten .
\end{equation}
The value is identical at both primes and stable under enlarging the sample set.
Twelve of the fourteen basis elements are graph invariants $I_{10,1},\ldots,
I_{10,12}$; the remaining two are the products $I_{4,1}I_{6,1}$ and
$I_{4,1}I_{6,2}$, which are the only products of lower-degree invariants
available at degree ten, since the ring begins at degree four and $4+6$ is the
sole splitting.

%=============================================================================
\section{Product, reachable and quotient decompositions}
\label{sec:decomp}

\subsection{Three decompositions, two of which look identical}

Three different splittings of the same fourteen-dimensional space are in play:
\begin{align}
  \dimAten &= \dimBten + 2  &&\text{published span } B_{10} + \text{complement},\\
  \dimAten &= \dimGten + \dimPten &&\text{graph generators } G_{10} + \text{products } P_{10},\\
  \dimAten &= \dimDten + \dimQten &&\text{reachable } D_{10} + \text{quotient } Q_{10}.
\end{align}
The first two both read $\dimBten + \dimPten$, which invites the identification
``the published candidates span the non-product part''.

\subsection{That identification is false}

Computed exactly at $\nIncidencePrimes$ primes:
\begin{equation}
  \dim(B_{10}\cap P_{10}) = \dimBcapP,
  \qquad
  \dim(B_{10} + P_{10}) = \dimBplusP \ne \dimAten .
\end{equation}
The published span therefore contains exactly one product direction and misses
one atlas direction. Its non-product content is $\publishedPrimitiveContent$, not
$\dimBten$, and the numerical agreement of the two twelves is a coincidence.

This is exact over the rationals, not only modulo a prime: the published
coordinates were lifted to $\mathbb{Q}$ by Chinese remaindering across
$\bpFitPrimes$ primes with rational reconstruction, validated at a held-out
prime~$\bpHoldoutPrime$, giving
$\dim_{\mathbb{Q}}(B_{10}\cap P_{10}) \bpCapRelation \bpCapQ$. The
intersection is one-dimensional and we can write its generator down: with
$I^{(k)}$ the published structures in the numbering of~\cite{Cederwall:2025invariants},
\begin{equation}
  - 55\,I^{(1)} - 252\,I^{(2)} + 4200\,I^{(5)} + 5040\,I^{(6)} \;=\; - 1040\,I^{(4)}_{1}I^{(6)}_{1}

  \label{eq:bpidentity}
\end{equation}
so a specific integer combination of four of the twelve published structures is
a multiple of a pure product. We verified~\eqref{eq:bpidentity} at a prime
$\bpFreshPrime$ and on $\bpFreshSamples$ five-forms that took no part in its
construction.

The decomposition that \emph{is} structural is
\begin{equation}
  P_{10}\cap G_{10} = 0,\quad P_{10}+G_{10} = \dimAten
  \qquad\Longrightarrow\qquad
  A_{10} = G_{10}\oplus P_{10}.
\end{equation}
Furthermore $P_{10}\subset D_{10}$: products contribute nothing to the quotient
$Q_{10}$. The complete pairwise incidence data are in
table~\ref{tab:incidence}.

\begin{table}[t]
\centering
\begin{tabular}{lrrrrl}
\toprule
pair $(X,Y)$ & $\dim X$ & $\dim Y$ & $\dim(X\cap Y)$ & $\dim(X+Y)$ & containment \\
\midrule
$A_{10}$, $B_{10}$ & 14 & 12 & 12 & 14 & $B10 \subset A10$ \\
$A_{10}$, $D_{10}$ & 14 & 11 & 11 & 14 & $D10 \subset A10$ \\
$A_{10}$, $G_{10}$ & 14 & 12 & 12 & 14 & $G10 \subset A10$ \\
$A_{10}$, $P_{10}$ & 14 & 2 & 2 & 14 & $P10 \subset A10$ \\
$B_{10}$, $G_{10}$ & 12 & 12 & 10 & 14 & neither \\
$D_{10}$, $B_{10}$ & 11 & 12 & 9 & 14 & neither \\
$D_{10}$, $G_{10}$ & 11 & 12 & 9 & 14 & neither \\
$P_{10}$, $B_{10}$ & 2 & 12 & 1 & 13 & neither \\
$P_{10}$, $D_{10}$ & 2 & 11 & 2 & 11 & $P10 \subset D10$ \\
$P_{10}$, $G_{10}$ & 2 & 12 & 0 & 14 & complementary \\
\bottomrule
\end{tabular}
\caption{Complete pairwise incidence of the degree-ten subspaces, identical at both primes. The two rows that matter most are $(P_{10},B_{10})$, whose intersection is nonzero, and $(P_{10},G_{10})$, whose intersection is zero.}
\label{tab:incidence}
\end{table}


\begin{figure}[t]
\centering
\includegraphics[width=0.92\textwidth]{figures/invariant_space_diagram.pdf}
\caption{The degree-ten space and its subspaces, drawn from the incidence
certificate. The graph generators and the products are complementary; the
published span is a third subspace which is not complementary to the products.}
\label{fig:spaces}
\end{figure}

\subsection{The quotient $Q_{10}$}

The stress-flow reachable subspace has $\dim D_{10} = \dimDten$, so
\begin{equation}
  \dim Q_{10} = \dim A_{10}/D_{10} = \dimQten .
\end{equation}
This is the concrete statement that the flow construction does not reach the
whole degree-ten space: there are $\dimQten$ independent directions of local
degree-ten structure that it cannot generate. Explicit compact representatives
are given in section~\ref{sec:pureN}.

\subsection{The quotient over the rationals, not only modulo $p$}
\label{sec:czero}

Every dimension above is a rank computed over $\mathbb{F}_p$, and for most of
them that costs nothing: a space spanned by exactly $k$ explicit invariants has
$\dim_{\mathbb{Q}} \le k$ for free, so a modular rank of $k$ pins the
dimension over $\mathbb{Q}$ with no prime excluded. This covers $A_{10}$,
$P_{10}$, $G_{10}$ and $B_{10}$.

$D_{10}$ is the exception, and it is the one that matters. It is not given by a
spanning set; it is assembled by a closure that admits a generated flow direction
when that direction raises the rank \emph{modulo $p$}. A direction rejected mod
$p$ still lies in $D_{10}$ and could be independent over $\mathbb{Q}$. So
$\dimDten$ bounded $\dim_{\mathbb{Q}} D_{10}$ only from below, and
$\dim_{\mathbb{Q}} Q_{10} \le \dimQten$ was all that followed.

Running the same fixed point in exact rational arithmetic removes the modular
admission test, and with it the gap. The obstacle is that the flow targets are
not all integral: of the targets at degrees four through ten,
$\czIntegerRows$ have identical integer coordinates at every prime and
$\czLiftedRows$ are reductions of rationals. Those were lifted by CRT across
$\czFitPrimes$ primes followed by rational reconstruction, and every
reconstructed entry was then required to reduce correctly at a held-out
prime~$\czHoldoutPrime$ that took no part in the lift. No lift failed and no
held-out reduction disagreed.

The closure then reaches its fixed point in $\czSweeps$ sweeps under exact
rational arithmetic, giving
\begin{equation}
  \dim_{\mathbb{Q}} D_{10} = \dimDtenQ ,
  \qquad
  \dim_{\mathbb{Q}} Q_{10} \czQtenRelation \dimQtenQ .
  \label{eq:czero}
\end{equation}
This is an equality and not a bound, because at a fixed point reached over
$\mathbb{Q}$ no generated direction raises the rank over $\mathbb{Q}$: the
reached span \emph{is} $D_{10}$. No separate upper-bound argument is needed --- the
upper bound was an artefact of computing in $\mathbb{F}_p$. As an independently
checkable witness the certificate records an explicit
$\czMinorSize \times \czMinorSize$ minor of the reached span whose integer
determinant, computed by fraction-free elimination, is non-zero.

What this does \emph{not} do is enlarge the generator set. The statement is that
the flow image generated by the targets recorded in the stress-flow certificates
has rank exactly $\dimDtenQ$ over $\mathbb{Q}$; whether that target list is
itself complete is a separate question and is unchanged.

\subsection{How many directions are needed, in any basis}
\label{sec:cardinality}

Everything above is a statement about dimensions computed in a particular graph
basis, and it is worth isolating the one consequence that does not depend on that
choice.

\begin{proposition}
\label{prop:cardinality}
Let $A$ be the degree-$d$ atlas, $D\subseteq A$ the reachable subspace, and
$\pi : A \to Q = A/D$ the quotient map. If a finite set $S\subseteq A$ closes the
degree, that is $D + \mathrm{span}(S) = A$, then $|S| \ge \dim Q$.
\end{proposition}

\begin{proof}
Applying $\pi$ to $D+\mathrm{span}(S)=A$ and using $\pi(D)=0$ and the
surjectivity of $\pi$ gives $\mathrm{span}(\pi(S)) = Q$. A span of $|S|$ vectors
has dimension at most $|S|$, so $\dim Q \le |S|$.
\end{proof}

The statement is basis-free: $A$, $D$ and $Q$ are defined without reference to
coordinates, and a change of basis of $A$ induces an isomorphism of $Q$ and
leaves $\dim Q$ unchanged. Since $\dim Q_{10} = \dimQten$ and the exhibited
completion has exactly $\dimQten$ elements, that completion is of \emph{minimum
cardinality in any basis} --- no relabelling or linear redefinition of the atlas
can close degree ten with fewer.

We are deliberately narrow about what this covers. It bounds the \emph{number} of
directions required. It says nothing about \emph{which} directions those are, and
it is not the stronger statement that no single element of the exhibited set may
be dropped. That removal-minimality is verified only in the fixed basis and under
relabellings of it, and is listed among the limitations of
section~\ref{sec:limitations}.

%=============================================================================
\section{Compact bases for $Q_{10}$}
\label{sec:pureN}

A basis is only useful if its elements can be written down. We give a
presentation in which each $Q_{10}$ representative is a single indexed
contraction, avoiding a notational ambiguity present in the source expressions.
The formulas, their serialisation, and the renderer/parser round trip are in
appendix~\ref{app:formulas}.

We do \emph{not} call this basis canonical. It is a choice, certified to span
$Q_{10}$ and validated on holdout samples; no universal canonicality is claimed
or proved.

%=============================================================================
\section{Formula-independent reverse recovery}
\label{sec:reverse}

To check that the quotient is not an artefact of how the invariants were
written, $Q_{10}$ is recovered a second time by a search that does not import the
published-formula evaluators at all. The search is organised into sectors; most
are capped, one is enumerated fully. It recovers rank three and its span agrees
with both the graph-derived and the published-candidate quotient spans.

We stress the scope: this is a \emph{bounded} search over a declared class. It is
not an exhaustive enumeration, and nothing in this paper depends on it being one.
The reverse recovery is corroboration, not the primary derivation.

%=============================================================================
\section{Independent spinor-variable calculation}
\label{sec:spinor}

The same invariant ring can be built in spinor variables. Ten dimensions admits Majorana--Weyl spinors in signatures with $s-t \equiv 0$
mod $8$~\cite{Kugo:1982bn}, which covers both $(1,9)$ and $(5,5)$. A chiral
spinor of $\mathrm{Spin}(10)$ has dimension $\spinorDim$, and
\begin{equation}
  \mathrm{Sym}^2(\spinorDim) = \dimSymSpinor = 10 \oplus \dimGammaTraceless,
\end{equation}
the $\dimGammaTraceless$ being the gamma-traceless part --- which as a module is
exactly the self-dual five-form. An independently written implementation
enumerates invariants in these variables.

\subsection{The frames are different, and the difference is not what was thought}
\label{sec:signature}

That implementation realises the ten vector gamma matrices as five wedge and five
contraction operators on $\Lambda^\bullet W$ with $\dim W = 5$. The metric this
induces follows from the exterior algebra alone,
\begin{equation}
  \{\Gamma^\mu,\Gamma^\nu\} = 2\eta^{\mu\nu},
  \qquad
  \eta = \tfrac12\begin{pmatrix} 0 & \mathbb{1}_5\\ \mathbb{1}_5 & 0\end{pmatrix},
\end{equation}
whose real signature is $(5,5)$ --- \emph{split}, not Euclidean. This matters
because in split signature, exactly as in Lorentzian,
$\star^2 = (-1)^{25}(-1) = +1$ on five-forms, so real self-dual five-forms exist.
A null frame of this kind cannot be Euclidean: Euclidean signature has no
nonzero isotropic vectors, whereas every basis vector here is null.

What does remain true is weaker and precise. Over $\mathbb{R}$, $(5,5)$ and
$(1,9)$ are inequivalent real forms of $\mathfrak{so}(10,\mathbb{C})$ and no real
matrix relates them. But both metrics have discriminant $-1$ modulo squares, so
over $\mathbb{C}$ --- and over $\mathbb{F}_p$, where signature does not exist and
a nondegenerate form is classified by its discriminant alone --- they are
congruent, and the transition is exact.

Two consequences, of different strengths, are worth separating.
\emph{Invariant dimensions are directly comparable} between the two
implementations with no complexification caveat at all: a real form is Zariski
dense in its complexification, so the graded dimensions of the invariant ring do
not depend on which real form is used. \emph{Component values, however, require
transport}, and we carry that out over $\mathbb{F}_p$.

%=============================================================================
\section{The spinor--tensor bridge}
\label{sec:bridge}

\subsection{The map}

With the frame transition $L$ constructed explicitly (rather than assumed to
exist), the bridge is
\begin{equation}
  S_{ab} \;=\; \frac{1}{5!}\,F_{\mu_1\ldots\mu_5}\,
                (\Gamma^{\mu_1\ldots\mu_5})_{ab}.
  \label{eq:bridge}
\end{equation}
Every convention entering~\eqref{eq:bridge} is pinned as a named constant:
signature, epsilon, Hodge star, self-duality sign, Clifford normalisation,
chirality, charge conjugation, antisymmetrisation weight and coefficient field.
Appendix~\ref{app:conventions} lists them.

\begin{figure}[t]
\centering
\includegraphics[width=0.92\textwidth]{figures/bridge_diagram.pdf}
\caption{The bridge. Both the kernel and the image are identified by two-way span
equality rather than by dimension counting, at each prime.}
\label{fig:bridge}
\end{figure}

\subsection{What is verified}

At each prime, and by two-way span containment rather than by dimension
counting:
\begin{itemize}
\item the kernel of~\eqref{eq:bridge} \emph{equals} the anti-self-dual
      $\dimAntiSelfdual$;
\item the image \emph{equals} the gamma-traceless $\dimGammaTraceless$;
\item the forward map has rank $\bridgeForwardRank$, not $\dimSymSpinor$;
\item the constructed left inverse satisfies
      $\mathrm{inverse}\circ\mathrm{forward} = P_{\text{self-dual}}$ exactly ---
      the identity on self-dual forms, and correctly \emph{not} the identity in
      general, since the anti-self-dual part is genuinely destroyed;
\item the map is equivariant, with the character solved for rather than assumed,
      under both a $GL(5)$ subgroup and products of Clifford reflections. The
      latter generate the full orthogonal group by Cartan--Dieudonné.
\end{itemize}
The bridge contains no floating-point step and no tolerance.

One convention was \emph{determined} rather than chosen. The spinor
implementation does not document whether its symmetric $\sigma^\mu_{ab}$ carries
upper or lower spinor indices. Of the eight candidate placements, exactly one
satisfies the equivariance requirement on every component with a single scalar;
that one was adopted. This is a reconstruction of an undocumented convention, and
is flagged as such.

%=============================================================================
\section{Common-sample comparison}
\label{sec:comparison}

A registry of $\nCommonSamples$ self-dual five-forms is fixed: $\nSampleSparse$
sparse, $\nSampleStructured$ structured, $\nSampleGeneric$ generic and
$\nSampleHoldout$ holdout. Every sample is hashed, and its self-duality and the
gamma-tracelessness of its image are verified before use. Both descriptions are
evaluated at the same five-forms in the same field.

\begin{table}[t]
\centering
\begin{tabular}{rrrlll}
\toprule
degree & trace rank & spinor rank & spans equal & holdout validated & spinor stopping \\
\midrule
4 & 1 & 1 & yes & yes & draw budget exhausted \\
6 & 2 & 2 & yes & yes & rank saturation (patience exhausted) \\
8 & 7 & 6 & NO & NO & rank saturation (patience exhausted) \\
10 & 14 & 14 & yes & yes & rank saturation (patience exhausted) \\
\bottomrule
\end{tabular}
\caption{Trace and spinor evaluation ranks on the common sample registry, exact over $\mathbb{F}_p$. Span equality is established by two-way containment; the change-of-basis map is fitted on the fitting samples and validated on the holdout samples. The spinor stopping column records that the spinor-side candidate stream is sampled to rank saturation, which is an observation and not a claim of exhaustive enumeration.
}
\label{tab:comparison}
\end{table}


\begin{figure}[t]
\centering
\includegraphics[width=0.62\textwidth]{figures/spinor_trace_rank_comparison.pdf}
\caption{Evaluation ranks on the common sample registry. The degree-eight bars
differ; that difference is a property of the spinor-side candidate family, not of
the correspondence, and is discussed in the text.}
\label{fig:comparison}
\end{figure}

Where the spans agree, the agreement is established by two-way containment and
the change-of-basis matrix is fitted on the fitting samples only and then
validated on the holdout samples. We do not report a degree as matched on the
strength of a dimension coincidence.

\paragraph{Degrees four, six and ten agree; degree eight does not, and that is a
result.} At degrees four, six and ten the two evaluation spans are equal, with
the change-of-basis map validated on holdout samples that took no part in the
fit. At degree ten in particular both sides reach $\traceRankDegTen$ and the
spans coincide, which is the cross-validation this comparison was built to test.

At degree eight the spinor-side port-graph stream reaches
$\spinorRankDegEight$ where the tensor side reaches $\traceRankDegEight$, and the
containment is strict: every spinor-side invariant lies in the tensor-side span,
but one tensor-side direction is not reached. This is not a defect of the bridge,
which is verified to be an isomorphism, nor of the sampling, which had saturated.
It is a property of the port-graph ansatz: single-graph contractions of the
invariant tensor miss a degree-eight direction.

This is not an inference. Running the spinor enumeration with its
stopping rule disabled, so that it terminates on candidate exhaustion rather than
on an externally supplied target, reaches rank seven at degree eight --- but only
when a separate family of structured tensor-word candidates is included alongside
the port graphs. Our exact modular computation re-implements the port graphs and
not that family, and lands on six. The two numbers are consistent, and together
they locate the missing direction precisely: it is not reachable by any
single-graph contraction.

This can be made sharper than ``the two numbers are consistent''. Re-running the
degree-eight comparison with the \emph{full} spinor candidate family --- port
graphs, products of lower degrees, the archive's structured candidates, and
tensor words, between $\dEightSpinorRowsMin$ and $\dEightSpinorRowsMax$ rows
depending on the prime --- gives
\begin{equation}
  \text{trace rank} = \text{spinor rank} = \text{union rank}
  = \dEightUnionRank
\end{equation}
at all $\dEightPrimes$ primes tested, with two-way containment and a
change-of-basis map validated on holdout samples. The union rank is the
load-bearing number: two \emph{different} seven-dimensional subspaces would also
give equal ranks, and a union of eight; the union being $\dEightUnionRank$ is
what makes this span equality rather than dimension agreement.

Dropping each spinor family in turn and re-ranking then isolates the
responsible one. Removing the port graphs, the products or the structured
candidates leaves the rank at $\dEightUnionRank$; removing the
\dEightIndispensable{} candidates drops it to $\dEightRankWithoutWords$. So the
direction that a port-graph-only family misses is not reachable by any
single-graph contraction of the invariant tensor, and the structured family does
not supply it either. That is a positive statement about where the invariant
lives, not merely an observation that two counts differ.

The consequence for reading table~\ref{tab:comparison} is that a ``no'' in the
degree-eight row is a statement about the candidate family, not about the
correspondence between the two descriptions.

\subsection{A degree-by-degree structural correspondence}

Both descriptions split each graded piece into non-product generators plus
products of lower-degree invariants. Those splits are computed independently, in
different variables and in different arithmetic. Where the spinor scan reaches a
terminal state they agree exactly:

\begin{center}
\begin{tabular}{rll}
\toprule
degree & tensor: graph $+$ product & spinor: new $+$ product \\
\midrule
4 & $1 + 0$ & $1 + 0$ \\
6 & $2 + 0$ & $2 + 0$ \\
8 & $6 + 1$ & $6 + 1$ \\
10 & $\dimGten + \dimPten$ & not reached by that scan \\
\bottomrule
\end{tabular}
\end{center}

This is the degree-by-degree form of $A_{10} = G_{10}\oplus P_{10}$. It is worth
being explicit about what it does not say: the published-candidate span $B_{10}$
is a \emph{different} twelve-dimensional subspace and is not a product
complement. The correspondence is with the graph generators, not with the
published structures --- which is precisely the distinction of
section~\ref{sec:decomp}.

%=============================================================================
\section{Exact Jacobian and the generic functional dimension}
\label{sec:jacobian}

\subsection{Why the previous computation was fragile}

The generic functional dimension $\genericFunctionalDim$
of~\eqref{eq:eightyone} is analytic. The previously available computational check
was a float64 central-difference Jacobian followed by a rank tolerance. That
procedure is genuinely fragile: two archived runs of the same code, differing
only in the random point, reported ranks of $35$ and $81$. The reason is not a
bug. In one run a large fraction of the candidate rows evaluated to numerical
zero at that point; row-normalising such rows rescales rounding noise into unit
vectors and manufactures rank.

Table~\ref{tab:jacruns} in appendix~\ref{app:jacobian} re-analyses both archived
Jacobians under an explicit rule: a row at or below the declared noise floor is
never normalised. The two runs separate cleanly. One has no rows at the noise
floor, a rank that is stable across six orders of magnitude of tolerance, and a
spectral gap of about $2\times 10^{7}$ at index $\genericFunctionalDim$; it is a
nondegenerate generic sample. The other has $48$ of $83$ rows at the noise floor,
no spectral gap at all, and an honest rank of $35$ --- and if the normalisation
rule is broken it reports $83$, which is the fabricated number. We keep the
degenerate example in the paper deliberately: it is the evidence that genericity
checks are necessary.

A full matrix over seeds, scales and step sizes was not run. A single
finite-difference Jacobian over these candidates costs more than ten minutes even
with parallel evaluation, so such a matrix was out of reach here; this is stated
rather than papered over. It matters less than it would otherwise, because the
exact modular computation below supersedes the floating-point evidence entirely,
and the archived pair already exhibits both behaviours.

\subsection{The exact replacement}

A port-graph invariant is multilinear in its factors of $F$, so its derivative
can be taken analytically by amputation,
\begin{equation}
  \frac{\partial I}{\partial F^{(k)}_{ab}}
  = \text{the same contraction with the $k$-th $F$ removed},
  \qquad
  \frac{\partial I}{\partial c_r} = \Big\langle \sum_k A_k,\; B_r \Big\rangle .
\end{equation}
No step size and no tolerance enter. The implementation is validated by Euler's
identity $\sum_r c_r\,\partial I/\partial c_r = (\deg I)\,I$, which holds exactly
at every degree tested.

Because the gamma-traceless basis is computed over $\mathbb{Z}$ --- its entries
are in $\{-1,0,+1\}$ and it is annihilated exactly --- the Jacobian is the
reduction of a genuine integer matrix, and
\begin{equation}
  \mathrm{rank}_{\mathbb{F}_p}(J \bmod p) \;\le\; \mathrm{rank}_{\mathbb{Q}}(J)
\end{equation}
holds unconditionally. An exact modular rank is therefore a \emph{rigorous} lower
bound in characteristic zero, not a probabilistic one.

An intermediate version of this computation covered only the port-graph subset of
the archive's candidate selection ($\nArchivePortGraphs$ of
$\nArchiveCandidates$) and reached exact modular rank $\archiveJacobianRank$. The
remaining $\nArchiveStructured$ structured tensor-word candidates have since been
implemented in the same exact arithmetic, and the certificate below covers the
complete selection.

\subsection{The certificate}
\label{sec:cert81}

Every one of the $\exactJacScheduled$ scheduled candidates carries a terminal
status: $\exactJacEvaluated$ are \texttt{evaluated}, with
$\exactJacErrors$ evaluation errors and $\exactJacZeroRows$ zero rows. Euler's
identity passes for every evaluated candidate (\exactJacEulerPass). The resulting
Jacobian is $\exactJacRows \times \exactJacCols$ and has exact modular rank
\begin{equation}
  \mathrm{rank}_{\mathbb{F}_p} J \;=\; \exactJacRank ,
  \label{eq:exactrank}
\end{equation}
reproduced independently at each of $\exactJacPoints$ certified sample
\exactJacPointWord{} over the \exactJacPrimeWord{} $\exactJacPrimes$, with the
per-degree cumulative ranks
\begin{center}
\begin{tabular}{lccccc}
\toprule
degree & 4 & 6 & 8 & 10 & 12 \\
\midrule
cumulative rank & \cumRankDegFour & \cumRankDegSix & \cumRankDegEight
                & \cumRankDegTen & \cumRankDegTwelve \\
\bottomrule
\end{tabular}
\end{center}

The rank statement is additionally witnessed by an explicit
$\minorSize \times \minorSize$ minor whose determinant is nonzero
(\minorDetNonzero) modulo $\minorPrimes$. Because the matrix is an
integer reduction, a single prime with nonvanishing determinant already proves
the integer minor is nonzero; the determinant is computed by two independent
routines, which agree, and additional primes guard against an indexing error in
selecting the minor rather than against a probabilistic failure. Hence
\begin{equation}
  \mathrm{rank}_{\mathbb{Q}} J \;\ge\; \charZeroLowerBound
\end{equation}
unconditionally --- no height bound, Hadamard estimate or rational reconstruction
is involved.

Two consequences deserve stating plainly. First, $\exactJacRows$ candidates of
rank $\exactJacRank$ means the selected candidates carry functional
\emph{dependencies}; they are not an independent set, and nothing here says they
are. Second, this lower bound meets the analytic upper bound
$\dimSelfDualModule - \dimLorentzAlgebra = \genericFunctionalDim$
of~\eqref{eq:eightyone} exactly. The two halves together pin the generic
functional dimension, but only the lower half is established here, and only at
finitely many points.

\paragraph{Permitted claim.} The generic functional dimension
$\genericFunctionalDim$ follows analytically from the trivial generic stabiliser.
The computations here provide an independent, exact, tolerance-free lower bound
on the same quantity. They did not discover $\genericFunctionalDim$ and are not
claimed to prove it by themselves.

%=============================================================================
\section{Physical interpretation}
\label{sec:physics}

\HUMAN{PHYSICS INTERPRETATION REQUIRES COAUTHOR CONFIRMATION}

The motivation for this ring is the classification of local algebraic
interactions of a chiral four-form gauge field in ten dimensions. A non-linear
theory of such a field is specified by a Lagrangian, or by a flow equation, built
from Lorentz invariants of $F$. In four and six dimensions the available
invariants are so few that a general deformation can be written explicitly. In
ten dimensions~\eqref{eq:eightyone} makes that impossible, and progress has
proceeded through specific constructions --- in particular stress-tensor
flows~\cite{Hutomo:2025chiral}.

The concrete consequence of the degree-ten result is a statement about what those
constructions reach. The stress-flow reachable subspace at degree ten has
dimension $\dimDten$ inside a $\dimAten$-dimensional space. So there are
$\dimQten$ independent directions of degree-ten local structure that the flow
cannot generate, and we exhibit representatives of all of them. This is an
equality over the rationals, not a bound: section~\ref{sec:czero} establishes
$\dim_{\mathbb{Q}} D_{10} = \dimDtenQ$ exactly.

This changes what a theorist can do in a specific way. Previously, a candidate
degree-ten term could be compared against the published structures only, and if
it failed to match, nothing followed --- the published list was not known to be
complete or independent. Now any degree-ten invariant can be expanded in a
certified basis, its coefficients read off exactly, and its class in $Q_{10}$
computed. That is the difference between recognising a term and locating it.

We deliberately do \emph{not} claim: a Type IIB coefficient, a supersymmetric
completion, a causality statement, an amplitude, or an all-orders theorem. None
of these follows from the algebra computed here.

%=============================================================================
\section{Limitations}
\label{sec:limitations}

\begin{enumerate}
\item \textbf{Modular, not characteristic zero --- and not uniformly so.} Lower
      bounds transfer to characteristic zero unconditionally; upper bounds do
      not. Applied to the individual results this separates them, and the
      separation is worth stating rather than covering with a blanket caveat.

      \emph{Unexposed.} A space spanned by exactly $k$ explicit elements has
      $\dim_{\mathbb{Q}} \le k$ for free, so if its modular rank is also $k$ then
      its dimension is $k$ over $\mathbb{Q}$ and no prime can be bad. This
      covers $\dim A_{10} = \dimAten$, $\dim P_{10} = \dimPten$,
      $\dim G_{10} = \dimGten$ and $\dim B_{10} = \dimBten$, each spanned by
      exactly that many explicit invariants. It also covers the Jacobian bound
      of section~\ref{sec:cert81}, for the separate reason that the matrix is an
      integer reduction and the claim is one-sided.

      \emph{Settled by exact computation.} $D_{10}$ was the one dimension
      assembled by a modular admission test, which made $\dimDten$ a lower bound
      over $\mathbb{Q}$ and $\dim_{\mathbb{Q}} Q_{10} \le \dimQten$ the
      supportable statement. Section~\ref{sec:czero} closes this: the closure
      run in exact rational arithmetic gives
      $\dim_{\mathbb{Q}} D_{10} = \dimDtenQ$ and hence
      $\dim_{\mathbb{Q}} Q_{10} = \dimQtenQ$.

      \emph{Also settled.} $\dim_{\mathbb{Q}}(B_{10}\cap P_{10})
      \bpCapRelation \bpCapQ$ likewise. The published-candidate coordinates were
      lifted to $\mathbb{Q}$ by the same route --- CRT across $\bpFitPrimes$
      primes, rational reconstruction, validation at a held-out
      prime~$\bpHoldoutPrime$ --- and the intersection is generated by an
      explicit integer identity, given in section~\ref{sec:decomp}.
\item \textbf{The reverse search is bounded.} It is a search over a declared
      class with most sectors capped. It is not an exhaustive enumeration and is
      never described as one.
\item \textbf{No canonicality claim.} The compact $Q_{10}$ basis is a certified
      choice, not a canonical object.
\item \textbf{Implementation independence, not clean-room independence.} The two
      implementations were written separately and neither imports the other;
      that is weaker than clean-room separation and is not claimed to be it.
\item \textbf{A source notational ambiguity is avoided, not resolved.} One
      published expression admits more than one index reading. We report the
      consequences of each reading and show which of them change coordinates
      without changing the span. We do not claim to have settled the source's
      intent.
\item \textbf{The signature bridge is modular.} Component-level comparison is
      carried out over $\mathbb{F}_p$, where the frame transition exists. A real
      Lorentzian-to-split component identification does not exist and is not
      claimed.
\item \textbf{Equivariance is verified at sampled group elements}, exactly, and
      is not a symbolic identity for all elements.
\item \textbf{One spinor convention is reconstructed.} The index placement in the
      spinor implementation's $\sigma^\mu_{ab}$ was determined by requiring
      equivariance, not read from documentation.
\item \textbf{Degree twelve is not classified.} It appears here only as
      background and outlook.
\item \textbf{Minimality is bounded in cardinality only.}
      Proposition~\ref{prop:cardinality} shows no basis admits a smaller closing
      set. It does not show that the particular exhibited set is irredundant
      under an arbitrary linear change of basis; that is verified in the fixed
      graph basis and under relabellings of it, and no more.
\item \textbf{The exact rank is a lower bound at finitely many points.} It is
      unconditional in characteristic zero because the Jacobian is an integer
      reduction, but it is evaluated at specific sample points and does not by
      itself establish a statement about a generic point. The matching upper
      bound is analytic and taken from the literature.
\end{enumerate}

%=============================================================================
\section{Conclusions and outlook}
\label{sec:conclusions}

The degree-ten graded piece of the invariant ring of the ten-dimensional
self-dual five-form is now known exactly, together with its product, published,
graph-generator and stress-flow-reachable subspaces and all their intersections.
The quotient by the reachable subspace is $\dimQten$-dimensional and has explicit
representatives. A tempting identification between two decompositions that both
read $\dimBten+\dimPten$ turns out to be false, and the correct structural
statement is $A_{10} = G_{10}\oplus P_{10}$.

Separately, the spinor-variable description has been connected to the tensor
description by an explicit convention-controlled bridge, and the numerical
Jacobian rank computation has been replaced by an exact modular one that yields a
rigorous bound.

The natural next targets are degree twelve, where the reachable subspace is
expected to be a much smaller fraction of the whole, and a generating set for the
ring itself. Neither is attempted here.

%=============================================================================
\acknowledgments
\HUMAN{ACKNOWLEDGMENTS}
\HUMAN{FUNDING}

%=============================================================================
\appendix

\section{Ten-dimensional conventions}
\label{app:conventions}
Every convention below is pinned as a named constant in
\texttt{spinor\_trace\_bridge/src/sdbridge/conventions.py}. Downstream code
imports the constant and may not re-derive a sign locally; a change to that file
is a scientific change and triggers a full re-run.

\paragraph{Metric.} Signature $(-,+,\dots,+)$, index $0$ timelike. The tensor
implementation's own metric function returns exactly this, and a test asserts the
equality rather than trusting it.

\paragraph{Epsilon and Hodge star.} $\epsilon_{01\ldots 9}=+1$;
$(\star F)_I = \frac{1}{(d-p)!}\epsilon_I{}^J F_J$, implemented on sorted
multi-indices as $\pm$ the complementary component times the raising signs, so
that no $10^{10}$ epsilon tensor is ever built. The value of $\star^2$ is
\emph{computed} at runtime, not assumed, and asserted to be $+\starSquared$.

\paragraph{Self-duality.} $\star F = +F$; projector $(1+\star)/2$, with $1/2$
taken as the modular inverse of $2$. Self-dual and anti-self-dual subspaces are
each $\dimSelfdual$-dimensional.

\paragraph{Clifford algebra.} $\{\Gamma^\mu,\Gamma^\nu\}=2\eta^{\mu\nu}$.
Chirality: $\Lambda^{\mathrm{even}}W$ is the positive-chirality module. The
Chevalley pairing carries the reversal sign $(-1)^{p(p-1)/2}$, which is what
makes $\sigma^\mu_{ab}$ symmetric; the symmetry is asserted, not assumed.

\paragraph{Antisymmetrised five-gamma.}
$(\Gamma^{\mu_1\ldots\mu_5})_{ab} = \frac{1}{5!}\sum_{\pi}\mathrm{sgn}(\pi)\,
(\sigma\bar\sigma\sigma\bar\sigma\sigma)_{ab}$, unit weight on the identity
permutation.

\paragraph{Forward normalisation.} Equation~\eqref{eq:bridge} sums over all index
tuples with the $1/5!$; because both objects are totally antisymmetric this
equals the sum over sorted tuples with unit weight, which is what the code does.

\paragraph{Inverse.} No closed form is claimed. A left inverse on the self-dual
subspace is constructed by inverting the $\dimSelfdual\times\dimSelfdual$
restriction in an explicit pivot basis; its defining property
$\mathrm{inverse}\circ\mathrm{forward}=P_{\text{self-dual}}$ is checked exactly.

\paragraph{Coefficient field.} $\mathbb{F}_p$ for $p\in\{\bridgePrimes\}$. The
bridge contains no floating-point step and no tolerance parameter.

\paragraph{Spinor index placement.} Determined, not chosen: see
section~\ref{sec:bridge}. Recorded as an open confirmation item.


\paragraph{Formulation independence of the leading-degree statement.}
Section~\ref{sec:czero}'s closure bookkeeping uses the fact that the free
stress-tensor trace vanishes at quadratic order, so that $\operatorname{Tr}\tau$
first contributes at field degree four. The vanishing does not depend on the
improvement convention: for any
$T_{\mu\nu} \sim F_{\mu\ldots}F_\nu{}^{\ldots} - c\,\eta_{\mu\nu}\langle F,F\rangle$
the trace is $(1-cd)\langle F,F\rangle$, a multiple of the single quadratic
scalar $\langle F,F\rangle$; and $\langle F,F\rangle$ vanishes identically on
either eigenspace of the star, because $F\wedge F$ is a top form built from two
copies of an odd-degree form and $\langle F,\star F\rangle$ is proportional to
it. A covariant action for a chiral form requires auxiliary
structure~\cite{Pasti:1995tn,Pasti:1997gx}, but the argument above uses only
the quadratic scalar available to the free theory, so it is insensitive to which
such formulation is adopted. Confirmation that the intended conventions are
those used here is coauthor review item~G-10.


\section{Degree-ten formulas}
\label{app:formulas}
The twelve degree-ten graph generators $I_{10,1},\dots,I_{10,12}$ and the two
products $I_{4,1}I_{6,1}$, $I_{4,1}I_{6,2}$ are serialised in the repository as
canonical contraction graphs, together with a renderer that emits an indexed
tensor expression and a parser that reads it back. The renderer/parser round trip
is tested: rendering a graph and re-parsing the result returns the same canonical
graph for every atlas element.

Compact representatives of $Q_{10}$ are given in a presentation in which each
representative is a single indexed contraction. This presentation was chosen to
avoid a notational ambiguity present in one published expression, where the index
placement admits more than one reading. We report the consequences of each
reading --- in particular which readings change coordinates while leaving the
span unaltered --- and we do not claim to have settled the source's intent.

Machine-readable forms, with hashes, are in \texttt{results/}; the mapping from
each formula to its atlas and quotient coordinates is in
\texttt{docs/PUBLISHED\_DEGREE10\_BASIS\_MAP.md} and
\texttt{docs/DEGREE10\_FORMULA\_CERTIFICATES.md}.


\section{Basis and quotient certificates}
\label{app:certificates}
All certificates are JSON with recorded hashes.

\begin{itemize}
\item \texttt{results/intrinsic\_candidates/}\allowbreak
      \texttt{degree10\_space\_incidence.json} ---
      dimensions of $A_{10},B_{10},G_{10},P_{10},D_{10}$ and the complete
      pairwise intersection/sum/containment data at each prime.
\item \texttt{results/intrinsic\_candidates/}\allowbreak
      \texttt{degree10\_published\_}\allowbreak\texttt{product\_intersection.json}
      --- the published-span/product-subspace intersection, with the superseded
      earlier conclusion preserved alongside the correction so that the change is
      auditable.
\item \texttt{results/rank81/certificate.json} --- the exact analytic Jacobian:
      a terminal status for every scheduled candidate, per-degree block and
      cumulative ranks, pivot rows and columns, the Euler homogeneity check, and
      per-candidate runtime and peak RSS.
\item \texttt{results/rank81/minor81\_certificate.json} --- the explicit
      $\minorSize\times\minorSize$ minor, its row and column indices, and its
      determinant computed by two independent routines.
\item \texttt{spinor\_trace\_bridge/results/}\allowbreak
      \texttt{bridge\_validation.json} --- Clifford
      relation, frame congruence, kernel and image span equalities, round trip,
      scaling, and the solved equivariance characters, at each prime.
\item \texttt{verification/COMMON\_SAMPLE\_REGISTRY.json} --- every common sample
      with its family, components and hash.
\item \texttt{verification/spinor\_trace\_comparison.json} --- per-degree ranks,
      span containments and holdout validation.
\item \texttt{verification/}\allowbreak
      \texttt{spinor\_archive\_jacobian\_exact.json} --- the exact
      modular Jacobian rank, with the excluded structured candidates listed
      explicitly rather than dropped.
\item \texttt{verification/SPINOR\_JACOBIAN\_RUNS.json} --- the float64 run
      matrix with per-run row-norm statistics, tolerance sweeps and
      classification.
\end{itemize}


\section{Algorithms and correctness tests}
\label{app:algorithms}
\paragraph{Canonicalisation.} Contraction graphs are reduced to a canonical form
so that two labellings of the same contraction are identified. A regression test
pins a known hash collision on regular multigraphs, so that a future change to
the hash cannot silently merge distinct graphs.

\paragraph{Modular rank.} A streaming row-reduced sieve accepts rows one at a
time and reports whether each is independent of everything seen so far. Memory is
proportional to the rank, not to the candidate count.

\paragraph{Span comparison.} Two subspaces are compared by two-way containment,
never by dimension. A dimension match alone would leave open that two spaces of
equal dimension are different subspaces, which is exactly the error that the
$12=12$ coincidence of section~\ref{sec:decomp} would otherwise have caused.

\paragraph{Contraction scheduling.} Spinor-side contractions are executed as
pairwise steps in an order found by a greedy path search, with a reduction mod
$p$ after each step so no intermediate overflows. Both a memory budget and a
floating-point-operation budget are enforced; a candidate exceeding either is
recorded as skipped, with the reason, rather than being silently dropped or
allowed to stall the run.

\paragraph{Derivative validation.} The analytic amputated derivative is validated
against Euler's identity $\sum_r c_r \partial_r I = (\deg I) I$, exactly, at every
degree used. The batched evaluator is separately validated against a literal
transcription of the port-graph definition.

\paragraph{Test counts.} The tensor-side suite has 199 tests; the bridge suite has
49, run at both primes.


\section{Jacobian audit}
\label{app:jacobian}
\begin{table}[t]\centering
\textbf{[float64 Jacobian run matrix: artifact not present at build time]}
\caption{float64 Jacobian run matrix --- pending.}
\label{tab:jacruns}
\end{table}


\paragraph{Why a degenerate run happens.} At certain sample points a large
fraction of the candidate rows evaluate to numerical zero. Their surviving
entries are rounding noise. Row-normalising rescales that noise to unit length,
after which the singular spectrum shows no gap and the reported rank becomes an
artefact of the tolerance. The rule enforced here is that a row whose norm is at
or below the declared noise floor is never normalised; the withheld
normalisation, and the reason, are recorded in the run record.

\paragraph{Why the exact computation is not subject to this.} The analytic
amputated derivative involves no step size, and rank over $\mathbb{F}_p$ involves
no tolerance. Because the gamma-traceless basis is integral --- entries in
$\{-1,0,+1\}$, annihilated exactly over $\mathbb{Z}$ --- the Jacobian is the
reduction of an integer matrix, so its modular rank is a rigorous lower bound on
the rank in characteristic zero.

\paragraph{Scope.} The exact computation covers the complete archived candidate
selection: $\exactJacScheduled$ scheduled, $\exactJacEvaluated$ evaluated,
$\exactJacErrors$ evaluation errors, $\exactJacZeroRows$ zero rows. Every
candidate carries a terminal status with a value, a derivative row, an output
hash, a runtime and a peak-RSS figure, so a candidate that failed would be
visible rather than silently absent from a rank.

An intermediate version of this appendix quoted $\archiveJacobianRank$ for the
port-graph subset alone ($\nArchivePortGraphs$ of $\nArchiveCandidates$), with
the $\nArchiveStructured$ structured tensor-word candidates excluded rather than
assumed inert. Those candidates were subsequently implemented in the same exact
arithmetic; section~\ref{sec:cert81} reports the result for the full selection.

\paragraph{What the number of primes does and does not buy.} For the rank
statement the prime count is not what carries the argument. An integer matrix
whose determinant vanishes vanishes modulo every prime, so a single prime at
which the explicit $\minorSize\times\minorSize$ minor has non-vanishing
determinant already establishes that the integer minor is nonzero, and hence the
characteristic-zero bound. No height bound, Hadamard estimate or rational
reconstruction enters. Additional primes are run because they guard against an
indexing error in selecting the minor --- a different failure mode, and one a
single prime would not catch. This is in contrast to the subspace dimensions of
section~\ref{sec:decomp}, which are equalities rather than one-sided bounds and
where additional primes do reduce genuine risk.


\section{Reproduction}
\label{app:reproduction}
\begin{table}[t]
\centering
\begin{tabular}{ll}
\toprule
item & value \\
\midrule
Python & 3.13.5 \\
NumPy & 2.5.1 \\
numpy & 2.5.1 \\
pynauty & 2.8.8.1 \\
pytest & 9.1.1 \\
primes & 32749, 32719 \\
\bottomrule
\end{tabular}
\caption{Reproducibility environment.}
\label{tab:env}
\end{table}


\paragraph{Quick reproduction.} From a clean clone at the recorded commit, with
the locked dependencies installed:
\begin{verbatim}
python -m pytest                                    # tensor side, 199 tests
python -m pytest -c spinor_trace_bridge/pytest.ini  # bridge, 49 tests
python manuscript/scripts/make_numbers.py
python manuscript/scripts/make_tables.py
python manuscript/scripts/make_figures.py
\end{verbatim}

\paragraph{Full reproduction.} The comparison and Jacobian certificates are
regenerated by the scripts in \texttt{spinor\_trace\_bridge/scripts/}. Each
writes its output incrementally and skips work already recorded, so a long run
can be interrupted and resumed.

\paragraph{Third-party code.} The spinor enumeration archive is not
redistributed with this work; redistribution permission has not been granted.
The release candidate contains a manifest with per-file hashes, the expected
local path, and adapter instructions, so that a reader holding their own copy can
reproduce the archive-dependent results. Results that do not depend on the
archive reproduce without it.


\bibliographystyle{JHEP}
\bibliography{references}

\end{document}
